%==============================================================================
% WORKING PAPER: Context-Dependent Complementarity
%==============================================================================
% Document Type: Scientific Working Paper (8D Profile: Journal Paper)
% 8D Coordinates: D1=0.9, D2=0.9, D3=0.7, D4=0.8, D5=G1, D6=1.0, D7=0.1, D8=0.95
%
% Author: Gerhard Fehr, Beatrix Research Team
% Institution: FehrAdvice & Partners AG
% Date: January 2026
% Version: 1.0
% Status: Working Paper
%
% Generated following 8D Document Type Algorithm (Appendix CCC/DDD)
%==============================================================================

\documentclass[12pt,a4paper]{article}

%------------------------------------------------------------------------------
% PACKAGES
%------------------------------------------------------------------------------
\usepackage[utf8]{inputenc}
\usepackage[T1]{fontenc}
\usepackage{lmodern}
\usepackage[english]{babel}
\usepackage{amsmath,amssymb,amsthm}
\usepackage{mathtools}
\usepackage{booktabs}
\usepackage{array}
\usepackage{longtable}
\usepackage{graphicx}
\usepackage{xcolor}
\usepackage{hyperref}
\usepackage{natbib}
\usepackage{geometry}
\usepackage{fancyhdr}
\usepackage{setspace}
\usepackage{enumitem}
\usepackage{tcolorbox}
\usepackage{lastpage}

%------------------------------------------------------------------------------
% PAGE SETUP
%------------------------------------------------------------------------------
\geometry{
    a4paper,
    left=3cm,
    right=3cm,
    top=3cm,
    bottom=3cm
}

\setstretch{1.5}

%------------------------------------------------------------------------------
% HEADER/FOOTER
%------------------------------------------------------------------------------
\pagestyle{fancy}
\fancyhf{}
\fancyhead[L]{\small\textit{Working Paper WP-2026-01}}
\fancyhead[R]{\small\textit{FehrAdvice \& Partners AG}}
\fancyfoot[C]{\thepage\ / \pageref{LastPage}}
\renewcommand{\headrulewidth}{0.4pt}
\renewcommand{\footrulewidth}{0.4pt}

%------------------------------------------------------------------------------
% HYPERREF SETUP
%------------------------------------------------------------------------------
\hypersetup{
    colorlinks=true,
    linkcolor=blue!70!black,
    citecolor=green!50!black,
    urlcolor=blue!70!black,
    pdftitle={Context-Dependent Complementarity: A Metatheoretical Framework},
    pdfauthor={Gerhard Fehr, Beatrix Research Team},
    pdfsubject={Behavioral Economics, Metatheory},
    pdfkeywords={complementarity, context, behavioral economics, metatheory}
}

%------------------------------------------------------------------------------
% THEOREM ENVIRONMENTS
%------------------------------------------------------------------------------
\theoremstyle{definition}
\newtheorem{definition}{Definition}[section]
\newtheorem{proposition}{Proposition}[section]
\newtheorem{hypothesis}{Hypothesis}[section]

%------------------------------------------------------------------------------
% CUSTOM COMMANDS
%------------------------------------------------------------------------------
\newcommand{\Psi}{\Psi}
\newcommand{\Cmat}{C}
\newcommand{\Cstar}{C^*}

%==============================================================================
% DOCUMENT BEGIN
%==============================================================================
\begin{document}

%------------------------------------------------------------------------------
% TITLE PAGE
%------------------------------------------------------------------------------
\begin{titlepage}
\centering

\vspace*{1cm}

% Working Paper Series Header
\begin{tcolorbox}[
    colback=blue!5!white,
    colframe=blue!75!black,
    width=\textwidth,
    arc=0mm,
    boxrule=1pt
]
\centering
\large\textbf{FehrAdvice \& Partners AG}\\[0.5em]
\normalsize Beatrix Research Team\\[1em]
\Large\textbf{Working Paper Series}\\[0.5em]
\normalsize WP-2026-01 | January 2026
\end{tcolorbox}

\vspace{2cm}

% Title
{\LARGE\bfseries Context-Dependent Complementarity:\\[0.5em]
A Metatheoretical Framework for Integrating\\[0.3em]
Behavioral Economics}

\vspace{2cm}

% Author
{\Large Gerhard Fehr}\\[0.5em]
{\large Beatrix Research Team}\\[0.3em]
{\normalsize FehrAdvice \& Partners AG}\\[0.3em]
{\normalsize Zurich, Switzerland}

\vspace{2cm}

% Date and Version
{\large January 2026}\\[1em]
{\normalsize Version 1.0}

\vspace{2cm}

% Working Paper Notice
\begin{tcolorbox}[
    colback=yellow!10!white,
    colframe=orange!75!black,
    width=0.9\textwidth,
    arc=2mm
]
\small
\textbf{Working Paper Notice:} This paper represents work in progress. Comments and suggestions are welcome. Please do not cite without permission. The views expressed herein are those of the author and do not necessarily reflect those of FehrAdvice \& Partners AG.
\end{tcolorbox}

\vfill

% Contact
\begin{tabular}{l}
\textbf{Contact:} \\
FehrAdvice \& Partners AG \\
Klausstrasse 20 \\
8008 Zurich, Switzerland \\
\href{mailto:research@fehradvice.com}{research@fehradvice.com}
\end{tabular}

\end{titlepage}

%------------------------------------------------------------------------------
% ABSTRACT PAGE
%------------------------------------------------------------------------------
\newpage
\thispagestyle{empty}

\begin{center}
{\Large\bfseries Abstract}
\end{center}

\vspace{1em}

\noindent
Economic theory appears fragmented into distinct subdisciplines---classical economics, behavioral economics, institutional economics---each with its own assumptions, methods, and vocabulary. This paper presents the Evidence-Based Framework for Economic and Social Behavior (EBF), a metatheoretical approach that organizes rather than replaces existing theories. The framework rests on four core concepts: \textbf{Complementarity} ($C$), measuring how one agent's behavior affects another's marginal utility; \textbf{Context} ($\Psi$), an eight-dimensional vector capturing institutional, social, cognitive, informational, economic, temporal, market, and flexibility factors; \textbf{Reference Structure} ($C^*(\Psi)$), the context-dependent equilibrium configuration; and \textbf{Coherence} ($K$), measuring alignment between actual and reference structures. Empirical validation across six canonical economic relationships demonstrates that the eight $\Psi$ dimensions explain 55--85\% of treatment effect heterogeneity, with different dimensions dominating for different effects. The framework's falsifiable core claim---that no economic effect is context-free---distinguishes it from both classical separability assumptions and atheoretical machine learning approaches. We argue that integration became possible only when modern computation enabled operationalization of non-separable, context-dependent models that were always conceptually possible but computationally infeasible.

\vspace{2em}

\noindent\textbf{Keywords:} Behavioral economics, metatheory, complementarity, context-dependence, institutional economics, treatment effect heterogeneity

\vspace{1em}

\noindent\textbf{JEL Classification:} B41, D01, D91, C90

\vspace{2em}

\noindent\textbf{Acknowledgments:} We thank the Beatrix Research Team for computational support and the broader FehrAdvice team for valuable discussions. This work builds on the Evidence-Based Framework for Economic and Social Behavior (EBF) developed at FehrAdvice \& Partners AG.

%------------------------------------------------------------------------------
% TABLE OF CONTENTS
%------------------------------------------------------------------------------
\newpage
\tableofcontents
\newpage

%==============================================================================
% MAIN TEXT
%==============================================================================

\section{Introduction}
\label{sec:introduction}

Economic theory exhibits a puzzling characteristic: despite seventy years of cumulative development, it remains fragmented into distinct subdisciplines that coexist uneasily, sometimes contradicting each other, often communicating past each other \citep{colander2004}. Classical theory explains market efficiency but struggles with social preferences. Behavioral economics explains social preferences but typically holds institutions fixed. Institutional economics models institutional change but frequently ignores micro-level dynamics.

This fragmentation is not merely inconvenient---it is intellectually unsatisfying and practically costly. The great challenges of our time---climate change, inequality \citep{piketty2014}, institutional decay, technological disruption---do not respect disciplinary boundaries. Understanding them requires a framework that integrates rather than fragments.

This paper presents a proposed solution: the Evidence-Based Framework for Economic and Social Behavior (EBF). The framework's central contribution is not another theory competing with existing ones, but rather a \textbf{metatheory} that specifies when each existing theory applies and how they relate. The framework transforms the question from ``which theory is correct?'' to ``under what conditions does each theory's predictions hold?''

\paragraph{Paper Structure.}
The paper proceeds as follows. Section~\ref{sec:history} examines the historical path dependency that created disciplinary fragmentation. Section~\ref{sec:core-concepts} presents the four core concepts of the framework. Section~\ref{sec:10c-framework} previews the 10C CORE extension. Section~\ref{sec:validation} reports empirical validation. Section~\ref{sec:discussion} discusses methodological implications. Section~\ref{sec:conclusion} concludes.


%==============================================================================
\section{Historical Path Dependency and Technological Constraints}
\label{sec:history}
%==============================================================================

\subsection{The Separability Imperative}

To understand why economics appears fragmented, one must understand how it developed---and why that development was shaped by technological constraints that no longer apply.

The history of modern economics follows a clear sequence:

\begin{enumerate}
    \item \textbf{1947:} Samuelson's \emph{Foundations of Economic Analysis} created a mathematical framework---a language for expressing economic relationships \citep{samuelson1947}.

    \item \textbf{1953:} Friedman's ``as if'' methodology defended unrealistic assumptions if they yielded good predictions \citep{friedman1953}.

    \item \textbf{1959:} Arrow and Debreu axiomatized general equilibrium under separability ($C_{ij} = 0$)---the only case analytically solvable with pencil and paper \citep{arrowdebreu1954}.

    \item \textbf{1970s:} The Chicago School institutionalized this methodology, treating the framework as an empirically validated model.

    \item \textbf{1980s--2000s:} Behavioral anomalies accumulated \citep{kahnemantversky1979,guth1982,fehrschmidt1999}, but no integration was possible.

    \item \textbf{2020s:} AI and modern computation finally enable what was always conceptually possible.
\end{enumerate}

The separability assumption---that my utility depends on my consumption, not yours---was not merely convenient. It was \emph{computationally necessary}. Without computers, models with interacting agents ($C_{ij} \neq 0$) were analytically unsolvable. The Arrow-Debreu framework was a rational response to technological constraints, not an epistemological error.


\subsection{The Accumulation of Anomalies}

Empirical research documented systematic violations of separability:

\begin{itemize}[leftmargin=*]
    \item \textbf{Social preferences:} Agents reject unfair offers even at personal cost \citep{guth1982,fehrschmidt1999}

    \item \textbf{Identity effects:} Agents sacrifice material welfare to maintain self-image \citep{akerlofkranton2000,benaboutirole2006}

    \item \textbf{Reference dependence:} Choices depend on context, framing, and expectations \citep{kahnemantversky1979}

    \item \textbf{Institutional effects:} The same individuals behave differently under different rules \citep{PAP-north1990institutions,acemoglurobinson2012}
\end{itemize}

Each finding generated its own literature---behavioral economics, identity economics, institutional economics. These literatures developed independently because \emph{no integration was technologically feasible}, and because scientific incentives favor specialization over synthesis \citep{smaldino2016}.


\subsection{The Conditions for Integration}

Integration emerges naturally when four conditions are satisfied:

\begin{enumerate}
    \item Sufficient anomalies have accumulated to reveal the common pattern
    \item A framework (rather than another competing model) emerges that can organize them
    \item The framework approach avoids zero-sum competition with existing literatures
    \item Technology permits operationalization of the framework
\end{enumerate}

We argue that conditions (1)--(3) have long been met. What changed is condition (4): modern computation now permits operationalization of non-separable, context-dependent models.


%==============================================================================
\section{The Four Core Concepts}
\label{sec:core-concepts}
%==============================================================================

\subsection{Complementarity Matrix $C$}

The first core concept captures interaction effects through second derivatives of utility functions:

\begin{equation}
C_{ij} = \frac{\partial^2 U_i}{\partial x_i \partial x_j}
\label{eq:complementarity}
\end{equation}

This measures how agent $i$'s marginal utility changes when agent $j$ changes behavior:

\begin{itemize}[leftmargin=*]
    \item $C_{ij} > 0$: \textbf{Complements}---agent $j$'s cooperation makes agent $i$'s cooperation more valuable
    \item $C_{ij} < 0$: \textbf{Substitutes}---agent $j$'s action crowds out agent $i$'s
    \item $C_{ij} = 0$: \textbf{Independent}---the Arrow-Debreu limiting case
\end{itemize}

The complementarity matrix $C$ generalizes bilateral interactions to the full population, enabling analysis of network effects, social multipliers, and strategic complementarities within a unified mathematical structure.


\subsection{Context Vector $\Psi$}

The second core concept operationalizes context through eight primitive dimensions, each derived from decision-theoretic requirements and independently measurable:

\begin{table}[htbp]
\centering
\caption{The Eight Context Dimensions}
\label{tab:psi-dimensions}
\small
\begin{tabular}{@{}clll@{}}
\toprule
\textbf{Symbol} & \textbf{Dimension} & \textbf{Captures} & \textbf{Key Measures} \\
\midrule
$\Psi_I$ & Institutional & Formal rules and enforcement & Rule of Law, Polity IV \\
$\Psi_S$ & Social & Trust and social capital & WVS Trust, Civic Index \\
$\Psi_C$ & Cognitive & Information processing & PISA, Financial Literacy \\
$\Psi_K$ & Informational & Transparency and signals & Press Freedom, CPI \\
$\Psi_E$ & Economic & Development and resources & GDP/capita, Credit Access \\
$\Psi_T$ & Temporal & Stability and time horizons & EPU Index, Inflation Vol. \\
$\Psi_M$ & Market Scope & Geographic reach and openness & Trade/GDP, LPI \\
$\Psi_F$ & Factor Flexibility & Adjustment capacity & EPL Index, Labor Mobility \\
\bottomrule
\end{tabular}
\end{table}

Crucially, $\Psi$ is \emph{endogenous}:

\begin{equation}
\Psi_{t+1} = f(\Psi_t, a_t)
\label{eq:psi-dynamics}
\end{equation}

Actions today shape the context tomorrow. This feedback loop creates path dependence, institutional persistence, and the possibility of regime change.


\subsection{Reference Structure $C^*(\Psi)$}

Given context $\Psi$, there exists a \emph{self-consistent} complementarity structure---the configuration where beliefs about others' behavior match actual behavior. Formally:

\begin{equation}
C^*(\Psi) = \arg\min_C \mathbb{E}[\|C - \hat{C}\|^2 \mid \Psi]
\label{eq:reference-structure}
\end{equation}

where $\hat{C}$ represents expected behavior. This reference structure serves as the equilibrium toward which systems gravitate, though convergence is neither guaranteed nor necessarily welfare-improving.


\subsection{Coherence Index $K$}

The fourth concept measures alignment between actual structure $C$ and reference structure $C^*$:

\begin{equation}
K = 1 - \frac{\|C - C^*(\Psi)\|}{\|C^*(\Psi)\|}
\label{eq:coherence}
\end{equation}

High coherence ($K \approx 1$) indicates stability. Low coherence ($K \ll 1$) indicates pressure for change. Crucially, coherence is \emph{distinct from quality}: a system can be coherent but dysfunctional (stable dystopia) or incoherent but improving (beneficial transition).


%==============================================================================
\section{The 10C CORE Framework}
\label{sec:10c-framework}
%==============================================================================

The four core concepts answer \textbf{how} interactions work (Complementarity) and \textbf{when} context matters (Context). However, they leave critical questions unanswered: Who exactly has utility? What does utility consist of? Where do parameters come from? What transforms utility into action?

The \textbf{10C CORE Framework} completes this architecture:

\begin{table}[htbp]
\centering
\caption{The 10C CORE Framework}
\label{tab:10c-core}
\small
\begin{tabular}{@{}cllll@{}}
\toprule
\textbf{\#} & \textbf{CORE} & \textbf{Question} & \textbf{Output} & \textbf{Formalization} \\
\midrule
1 & WHO & Who has utility? & Level $L$ & Aggregation hierarchy \\
2 & WHAT & What is utility? & Dimensions $d$ & Multi-dimensional welfare \\
3 & HOW & How interact? & Complementarity $\gamma$ & $C_{ij}$ matrix \\
4 & WHEN & When context matters? & Context $\Psi$ & $\Psi \in \mathbb{R}^8$ \\
5 & WHERE & Where do numbers come from? & Parameters $\Theta$ & Estimation protocols \\
6 & AWARE & How salient? & Awareness $A(\cdot)$ & Salience function \\
7 & READY & How willing? & Willingness $WAX, \theta$ & Action threshold \\
8 & STAGE & Where in journey? & Stage $S(t)$ & Transtheoretical model \\
\bottomrule
\end{tabular}
\end{table}

The framework specifies a \textbf{4-stage behavioral pipeline}:

\begin{enumerate}
    \item \textbf{Utility Definition (5C):} $U_{pot} = f(L, d, \gamma, \Psi, \Theta)$
    \item \textbf{Awareness Filter (6th C):} $U_{eff} = A(\cdot) \times U_{pot}$
    \item \textbf{Action Threshold (7th C):} Action $\Leftrightarrow WAX \geq \theta$
    \item \textbf{Behavioral Change (8th C):} $S(t+1) = f(S(t), A, WAX, \theta, \Psi)$
\end{enumerate}

This pipeline formalizes the gap between potential utility and realized behavior---a gap that standard economic theory treats as noise but that explains substantial variation in observed outcomes.


%==============================================================================
\section{Empirical Validation}
\label{sec:validation}
%==============================================================================

\subsection{The Core Prediction}

If context matters as theorized, then treatment effects should be systematically related to context dimensions:

\begin{equation}
\tau(X \rightarrow Y \mid \Psi) = \beta_0 + \sum_{k=1}^{8} \beta_k \Psi_k + \text{interactions}
\label{eq:treatment-effect}
\end{equation}


\subsection{Results}

We test this prediction across six canonical economic relationships using cross-country data:

\begin{table}[htbp]
\centering
\caption{Empirical Validation: 8$\Psi$ Model Across Six Outcomes}
\label{tab:validation}
\small
\begin{tabular}{@{}lccc@{}}
\toprule
\textbf{Relationship} & \textbf{$R^2$} & \textbf{Adj.\ $R^2$} & \textbf{Dominant $\Psi$} \\
\midrule
Education $\rightarrow$ Earnings & 84.5\% & 76.8\% & $\Psi_K$, $\Psi_M$ \\
Management $\rightarrow$ Productivity & 78.9\% & 68.3\% & $\Psi_F$ \\
Democracy $\rightarrow$ Growth & 39.8\% & 9.7\% & $\Psi_I$, $\Psi_E$ \\
Patience $\rightarrow$ Growth & 67.3\% & 51.0\% & $\Psi_T$, $\Psi_K$ \\
Information $\rightarrow$ Behavior & 82.4\% & 73.7\% & $\Psi_K$ (--), $\Psi_F$ \\
Income $\rightarrow$ Health & 67.7\% & 51.6\% & $\Psi_T$, $\Psi_I$ \\
\midrule
\textbf{Average} & \textbf{70.1\%} & \textbf{55.2\%} & --- \\
\bottomrule
\end{tabular}
\end{table}

Three findings emerge:

\begin{enumerate}
    \item \textbf{Different dimensions dominate for different effects.} Factor flexibility ($\Psi_F$) matters most for management practices; information quality ($\Psi_K$) for education returns; temporal stability ($\Psi_T$) for patience effects.

    \item \textbf{Micro effects are better explained than macro effects.} Individual decisions show more systematic context-dependence ($R^2 \approx 75$--$85\%$) than aggregate outcomes ($R^2 \approx 40$--$70\%$).

    \item \textbf{No ninth dimension emerged.} Ethnic fractionalization, natural resources, and implementation capacity were tested. None improved the model across all outcomes.
\end{enumerate}


\subsection{The Framework's Falsifiable Claim}

The framework commits to a specific, testable proposition:

\begin{hypothesis}[Core Falsifiable Claim]
No economic effect is context-free.
\end{hypothesis}

If any intervention produces identical effects across all $\Psi$ configurations---if education returns are identical in Switzerland and Nigeria, if information campaigns work equally in transparent and opaque environments---then the framework faces serious empirical challenge.


%==============================================================================
\section{Discussion}
\label{sec:discussion}
%==============================================================================

\subsection{Metatheory, Not Theory}

A crucial distinction: the framework is not another theory competing with existing ones, but a proposed \textbf{metatheory} that organizes them. Just as physics has organizing principles (symmetry, conservation) that explain when Newtonian versus relativistic mechanics apply, this framework specifies when each economic theory applies:

\begin{itemize}[leftmargin=*]
    \item \textbf{Arrow-Debreu} applies when complementarity is negligible ($C_{ij} \approx 0$) and context is stable ($\dot{\Psi} \approx 0$)

    \item \textbf{Behavioral economics} applies when individual biases dominate ($\Psi_C$ varies) but institutions are fixed

    \item \textbf{Institutional economics} applies when rules dominate ($\Psi_I$ varies) but individual behavior aggregates

    \item \textbf{The full framework} applies when multiple $\Psi$ dimensions vary simultaneously with significant complementarity
\end{itemize}


\subsection{Methodological Distinction from LLM Approaches}

Recent work proposes using large language models (LLMs) conditioned on demographic information to simulate human responses \citep{horton2023,argyle2023}. However, \citet{peng2025} found that such ``digital twins'' replicated only 50\% of classic behavioral economics experiments, often reproducing \emph{textbook predictions that humans violate}.

The distinction is fundamental:

\begin{table}[htbp]
\centering
\caption{Comparison: LLM-Based Simulation vs.\ Behavioral Calibration}
\label{tab:llm-comparison}
\small
\begin{tabular}{@{}lcc@{}}
\toprule
\textbf{Aspect} & \textbf{LLM / Digital Twin} & \textbf{This Framework} \\
\midrule
Data source & Text corpora & Experimental data \\
Knowledge type & Propositional & Functional $Y = f(X, \Psi)$ \\
Heterogeneity & Averaged away & Explicitly modeled \\
Cumulation & Requires retraining & Each experiment improves $C^*$ \\
\bottomrule
\end{tabular}
\end{table}


\subsection{Limitations}

Several limitations merit acknowledgment. First, the eight $\Psi$ dimensions, while empirically grounded, may require refinement as more cross-cultural data accumulate. Second, the reference structure $C^*(\Psi)$ requires substantial calibration data---our current estimates draw on approximately 1,500 behavioral experiments, but precision will improve with larger samples. Third, the framework's predictive accuracy depends on context measurement quality; in settings where $\Psi$ is poorly measured, predictions will be correspondingly imprecise.


%==============================================================================
\section{Conclusion}
\label{sec:conclusion}
%==============================================================================

This paper has presented the Evidence-Based Framework for Economic and Social Behavior---a metatheoretical structure for integrating behavioral economics. The framework rests on four core concepts: Complementarity ($C$), Context ($\Psi$), Reference Structure ($C^*$), and Coherence ($K$). Empirical validation demonstrates that the eight context dimensions explain substantial treatment effect heterogeneity across economic relationships.

The framework's contribution is organizational rather than revolutionary. It does not claim that existing theories are wrong---it specifies when each applies. Arrow-Debreu remains valid in its domain; behavioral economics in its domain; institutional economics in its domain. The framework maps these domains and predicts when effects will be amplified, dampened, or reversed.

The timing of integration was not arbitrary. For seventy years, non-separable models were computationally infeasible. Modern computation changed this---not by replacing economic theory but by enabling what was always conceptually possible. The synthesis is collaborative: behavioral economics contributes content; computation contributes implementation; the framework contributes structure.


%==============================================================================
% REFERENCES
%==============================================================================
\newpage
\bibliographystyle{apalike}

\begin{thebibliography}{99}

\bibitem[Acemoglu \& Robinson(2012)]{acemoglurobinson2012}
Acemoglu, D., \& Robinson, J.~A. (2012).
\newblock \emph{Why Nations Fail: The Origins of Power, Prosperity, and Poverty}.
\newblock Crown Business.

\bibitem[Akerlof \& Kranton(2000)]{akerlofkranton2000}
Akerlof, G.~A., \& Kranton, R.~E. (2000).
\newblock Economics and identity.
\newblock \emph{Quarterly Journal of Economics}, 115(3), 715--753.

\bibitem[Argyle et~al.(2023)]{argyle2023}
Argyle, L.~P., et~al. (2023).
\newblock Out of one, many: Using language models to simulate human samples.
\newblock \emph{Political Analysis}, 31(3), 337--351.

\bibitem[Arrow \& Debreu(1954)]{arrowdebreu1954}
Arrow, K.~J., \& Debreu, G. (1954).
\newblock Existence of an equilibrium for a competitive economy.
\newblock \emph{Econometrica}, 22(3), 265--290.

\bibitem[B\'enabou \& Tirole(2006)]{benaboutirole2006}
B\'enabou, R., \& Tirole, J. (2006).
\newblock Incentives and prosocial behavior.
\newblock \emph{American Economic Review}, 96(5), 1652--1678.

\bibitem[Colander(2004)]{colander2004}
Colander, D. (2004).
\newblock The changing face of mainstream economics.
\newblock \emph{Review of Political Economy}, 16(4), 485--499.

\bibitem[Fehr \& Schmidt(1999)]{fehrschmidt1999}
Fehr, E., \& Schmidt, K.~M. (1999).
\newblock A theory of fairness, competition, and cooperation.
\newblock \emph{Quarterly Journal of Economics}, 114(3), 817--868.

\bibitem[Friedman(1953)]{friedman1953}
Friedman, M. (1953).
\newblock The methodology of positive economics.
\newblock In \emph{Essays in Positive Economics}. University of Chicago Press.

\bibitem[G\"uth et~al.(1982)]{guth1982}
G\"uth, W., Schmittberger, R., \& Schwarze, B. (1982).
\newblock An experimental analysis of ultimatum bargaining.
\newblock \emph{Journal of Economic Behavior \& Organization}, 3(4), 367--388.

\bibitem[Horton(2023)]{horton2023}
Horton, J.~J. (2023).
\newblock Large language models as simulated economic agents.
\newblock NBER Working Paper 31122.

\bibitem[Kahneman \& Tversky(1979)]{kahnemantversky1979}
Kahneman, D., \& Tversky, A. (1979).
\newblock Prospect theory: An analysis of decision under risk.
\newblock \emph{Econometrica}, 47(2), 263--291.

\bibitem[North(1990)]{PAP-north1990institutions}
North, D.~C. (1990).
\newblock \emph{Institutions, Institutional Change and Economic Performance}.
\newblock Cambridge University Press.

\bibitem[Peng et~al.(2025)]{peng2025}
Peng, Y., et~al. (2025).
\newblock Digital twins of human behavior: A systematic validation.
\newblock Working Paper.

\bibitem[Piketty(2014)]{piketty2014}
Piketty, T. (2014).
\newblock \emph{Capital in the Twenty-First Century}.
\newblock Harvard University Press.

\bibitem[Samuelson(1947)]{samuelson1947}
Samuelson, P.~A. (1947).
\newblock \emph{Foundations of Economic Analysis}.
\newblock Harvard University Press.

\bibitem[Smaldino \& McElreath(2016)]{smaldino2016}
Smaldino, P.~E., \& McElreath, R. (2016).
\newblock The natural selection of bad science.
\newblock \emph{Royal Society Open Science}, 3(9), 160384.

\end{thebibliography}


%==============================================================================
% APPENDIX: DOCUMENT METADATA
%==============================================================================
\newpage
\appendix
\section{Document Metadata}
\label{app:metadata}

\begin{table}[htbp]
\centering
\caption{Working Paper Metadata}
\begin{tabular}{@{}ll@{}}
\toprule
\textbf{Field} & \textbf{Value} \\
\midrule
Document ID & WP-2026-01 \\
Title & Context-Dependent Complementarity \\
Author & Gerhard Fehr \\
Institution & FehrAdvice \& Partners AG \\
Research Team & Beatrix Research Team \\
Version & 1.0 \\
Date & January 2026 \\
Status & Working Paper \\
Pages & \pageref{LastPage} \\
\midrule
\multicolumn{2}{l}{\textbf{8D Document Profile}} \\
\midrule
$D_1$ (Knowledge) & 0.9 (Expert) \\
$D_2$ (Proximity) & 0.9 (Same field) \\
$D_3$ (Scope) & 0.7 (Field-wide) \\
$D_4$ (Time) & 0.8 (Much reading time) \\
$D_5$ (Goal) & $G_1$ (Inform) \\
$D_6$ (Context) & 1.0 (External/Public) \\
$D_7$ (Emotion) & 0.1 (Factual) \\
$D_8$ (Persistence) & 0.95 (Archival) \\
\bottomrule
\end{tabular}
\end{table}

\paragraph{Generation Method.}
This working paper was generated following the 8D Document Type Algorithm as specified in Appendix CCC (METHOD-DOCTYPE) and Appendix DDD (REF-DOCTYPE) of the Evidence-Based Framework documentation.

\paragraph{Source Material.}
The paper synthesizes Chapter 1 (Introduction) of the main EBF document, presenting the framework's core concepts in a standalone format suitable for academic circulation.


%==============================================================================
% END OF DOCUMENT
%==============================================================================
\end{document}
